	\begin{center}
		\section{博弈论}
	\end{center}
	% 这一部分的部分内容参考了下面这个链接：
	% https://github.com/hh2048/XCPC/blob/484e9e0e6e4f127f187eb3df5bc419e9ec863795/02%20-%20%E6%89%93%E5%8D%B0%E7%A8%BF%E6%A8%A1%E6%9D%BF%E6%B1%87%E6%80%BB/11%20-%20%E5%8D%9A%E5%BC%88%E8%AE%BA.md?plain=1#L92

	\begin{multicols*}{2}

		\subsection{Nim 博弈}

		\textbf{规则}：桌上有 $n$ 堆石子，第 $i$ 堆有 $a_i \ge 0$ 颗。两人轮流操作，每回合选一堆\textbf{非空}的石子，从中拿走任意正整数颗（至少 $1$ 颗，不能跨堆）。\textbf{无法操作}（所有堆都空）的玩家落败。问先手能否取胜。

		\textbf{结论（Bouton 定理）}：先手胜当且仅当
		\[
			a_1 \oplus a_2 \oplus \dots \oplus a_n \;\neq\; 0.
		\]
		把这个异或值称为局面的 \textcolor{blue}{\textbf{Nim 和}}（Nim-sum）。Nim 和为零的局面是 \textcolor{blue}{\textbf{P 局面}}（Previous，前一手获胜，即\textbf{当前轮到的人必败}）；非零是 \textcolor{blue}{\textbf{N 局面}}（Next，当前轮到的人必胜）。

		\textbf{为什么是异或？三段闭合性}：

		\begin{enumerate}
			\item \textbf{终止 = P}：全 $0$ 局面 Nim 和为 $0$，确实无路可走、当前人输。
			\item \textbf{P 走完必入 N}（封死 P 出口）：从 Nim 和 $=0$ 出发，任何合法操作只动一堆且必使该堆严格变小，单堆变化必然改变异或——\textcolor{red}{\textbf{改一堆就一定让 Nim 和从 $0$ 变非 $0$}}。所以 P 局面下你怎么走都把回合白送给对手。
			\item \textbf{N 总能走回 P}（构造 N 通道）：从 Nim 和 $s \neq 0$ 出发，找 $s$ 的最高位 $k$。$s$ 在第 $k$ 位为 $1$ 意味着 $a_i$ 中至少有一个在第 $k$ 位为 $1$，挑这样的一堆 $a_j$，把它改为 $a_j \oplus s$。两个关键事实：(i) $a_j$ 第 $k$ 位被 $s$ 翻成 $0$、更高位不动，所以 $a_j \oplus s < a_j$，是\textbf{合法的"减一个正数"操作}；(ii) 操作后总异或变成 $s \oplus s = 0$。因此 N 局面总能一步把 P 局面"踢回"对手。
		\end{enumerate}

		三段闭合后再做归纳：终止 P 局面下手必败 $\Rightarrow$ 一切 P 局面下手必败（因为只能进入 N，再被对手踢回 P 直至终止）；一切 N 局面下手必胜（一步进入 P，把上述循环转嫁给对手）。

		\textbf{Nim 走法（构造性）}：判负只是判定，但 \textcolor{red}{\textbf{遇到 N 局面时必须知道"该动哪一堆、动到几"}}。算法照搬上面 (3)：

		\begin{minted}{cpp}
// 返回 (堆下标 j, 改后的新值)；保证 a_j 严格减小
pair<int, ll> nim_move(const vector<ll>& a) {
    ll s = 0;
    for (ll x : a) s ^= x;
    if (s == 0) return {-1, -1}; // P 局面，无必胜走法
    for (int j = 0; j < (int) a.size(); ++j) {
        ll t = a[j] ^ s;
        if (t < a[j]) return {j, t}; // 找到含最高位的那一堆
    }
    return {-1, -1}; // 不会到这
}
		\end{minted}

		注意"含 $s$ 最高位"的判定可以\textbf{不显式取最高位}：直接看 $a_j \oplus s < a_j$ 是否成立——等价于 $a_j$ 在 $s$ 的最高位上为 $1$。这个写法少踩位运算坑。

		\textbf{推广：Sprague–Grundy 定理}。对任意\textbf{公平组合游戏}（双方走法对称、有限步必终止、无随机），定义局面 $u$ 的 \textcolor{blue}{\textbf{SG 值}} $\operatorname{SG}(u) = \operatorname{mex}\{\operatorname{SG}(v) : u \to v\}$。则：

		\begin{itemize}
			\item 单局游戏 $u$ 是 P 局面 $\iff \operatorname{SG}(u) = 0$。
			\item 多个独立游戏\textbf{并行}进行（每回合任选一局走一步），整体 SG 值 $= \bigoplus_i \operatorname{SG}(u_i)$。
		\end{itemize}

		Nim 单堆的 SG 值就是堆里石子数（堆 $a$ 一步可达 $\{0, 1, \dots, a-1\}$，mex 即 $a$），所以多堆 Nim 的"异或为零判 P"正是 SG 定理在 Nim 上的特化。\textcolor{red}{\textbf{遇到看似复杂的博弈，先尝试拆成若干独立子游戏 + 求每个 SG + 异或起来}}——下一节的"三角形博弈"就是把"非标准的单游戏"通过换元映回标准 Nim 的最简案例。

		\subsubsection{三角形博弈（带约束的 Nim）}

		给定能构成非退化三角形的三整数 $a, b, c$（即满足 $a+b>c$、$b+c>a$、$a+c>b$）。两人轮流操作，每回合选一个数减去任意正整数；操作后若三数无法构成非退化三角形，则该玩家落败。先手能否获胜？

		\textbf{结论}：先手胜当且仅当
		\[
			(a-1) \oplus (b-1) \oplus (c-1) \;\neq\; 0.
		\]

		\textbf{核心观察——把"边长"换元为"石子堆"}：令 $x_i = s_i - 1$（$s_i$ 是边长）。每条边至少为 $1$，对应每堆石子至少为 $0$；"减去任意正整数"在两套语言下完全一致。终止局面 $(1,1,1)$ 对应 $(0,0,0)$，与标准 Nim 的零局面重合。\textcolor{red}{\textbf{换言之，剥掉三角约束之后，这就是三堆 Nim}}。剩下要论证的只是：三角约束\textbf{不会挡掉}从 N 局面走向 P 局面的那一步（沿用上一节的 Nim 走法）。

		\textbf{为什么三角约束不挡 Nim 走法}：从异或非零的 $(a,b,c)$ 出发，设 $s = (a-1)\oplus(b-1)\oplus(c-1)$，找到 $(a-1, b-1, c-1)$ 中含 $s$ 最高位的那一堆，标准 Nim 操作把它改成"它 $\oplus s$"——这是一次\textbf{严格变小}的修改。不妨设被改的是 $c-1$，新值 $c' - 1 = (c-1)\oplus s = (a-1)\oplus(b-1)$。此时 $|a - b| \le c' - 1 \le a + b - 2 < a + b$，新三角不等式自动成立（最长边假设依然成立时同理），所以"修复异或"的 Nim 走法\textbf{始终是合法着}。

		\textbf{反向也封死}：从异或为零的 P 局面出发，任何合法操作只动一堆，异或必然变成非零；对手回到 N 局面继续兑现获胜策略。

		\textbf{实战提示}：实现一行就够——读入三个数后输出 \texttt{((a-1)\^{}(b-1)\^{}(c-1)) ? "Win" : "Lose"}。范例题：\href{https://qoj.ac/problem/4545}{QOJ 4545 Triangle Game}（凯特和艾米丽科，$1 \le a,b,c \le 10^9$，$T \le 10^4$，$O(1)$ 一次输出）。\textbf{推广视角}：所有"看似有附加约束的 Nim"都值得先做这套换元——\textcolor{red}{\textbf{把约束的下界吸收到偏移里、把"零状态"对齐成 $(0,0,\dots)$、再核对"最优 Nim 走法是否被约束挡掉"，三步通吃}}；本题正是 $-1$ 偏移加"最长边可调"的最简版本。

		\subsubsection{Moore's Nim-K（每回合可同时动 K 堆）}

		\textbf{规则}：$n$ 堆石子，每回合\textbf{至多选 $K$ 堆}，每选中的堆都取走任意正整数颗（不同堆取数可以不同），\textbf{取到最后一颗者胜}。$K=1$ 即标准 Nim。

		\textbf{结论}：把每堆数写成二进制，记第 $i$ 位上 $1$ 的个数为 $\operatorname{ones}_i$。\textcolor{red}{\textbf{先手败 $\iff$ 每个 $\operatorname{ones}_i$ 都是 $(K+1)$ 的倍数}}。$K=1$ 退化成"每位 $\operatorname{ones}_i$ 都是偶数"，即 XOR $=0$。

		\textbf{直觉}：标准 Nim 的"模 $2$ 异或"是 $K=1$ 的特例，因为每回合只能动 $1$ 堆，每一位上的 $1$ 个数变化量 $\in \{-1, 0, +1\}$，对手可恢复奇偶。Nim-K 下变化量 $\in \{-K, \dots, K\}$（每选中堆贡献 $\pm 1$），对手可恢复"模 $K+1$"残差。

		\subsubsection{反 Nim（Anti-Nim / Misère Nim）}

		\textbf{规则}：与标准 Nim 完全相同，但\textbf{取到最后一颗者输}（拿到最后一颗的"出局"）。

		\textbf{结论}：\textcolor{red}{\textbf{先手胜 $\iff$ 满足下列其一}} ——
		\begin{itemize}
			\item 所有堆都 $\le 1$ 且 Nim 和 $=0$（等价于"全是 $1$ 且堆数为偶数"）；
			\item 至少一堆 $> 1$ 且 Nim 和 $\neq 0$。
		\end{itemize}

		\textbf{直觉}：当所有堆都 $\le 1$ 时游戏退化成"轮流取走单颗"，胜负由堆数奇偶决定（misère 下偶数堆先手胜）。当存在 $>1$ 的"大堆"时，先手始终能控制把局面引入"全 $1$ 局面，且堆数对自己有利"——具体策略类似标准 Nim，只是终态判定反过来。"$>1$ 堆 + Nim 和非零"是先手能掌控终态奇偶的标志。

		\subsubsection{阶梯 Nim}

		\textbf{规则}：$n$ 级台阶（编号 $1 \dots n$），第 $i$ 级有 $a_i$ 颗石子。每回合从某级 $i \ge 1$ 拿任意正数颗放到第 $i-1$ 级（第 $1$ 级的石子被放到地面 $0$ 级，地面石子不能再动），\textbf{取走最后一颗者胜}（即让对手面对"无可动石子"局面者胜）。

		\textbf{结论}：\textcolor{red}{\textbf{只看奇数级台阶上的石子做标准 Nim}}——先手败 $\iff a_1 \oplus a_3 \oplus a_5 \oplus \dots = 0$。

		\textbf{直觉——"偶数级是镜子"}：偶数级上的动作可以被对手\textbf{原样模仿}——你从第 $2k$ 级搬 $x$ 颗到 $2k-1$ 级，对手就从 $2k-1$ 级搬同样 $x$ 颗到 $2k-2$ 级，把你"白送"到偶数级的石子继续往下推，整体偶数级局面回到等价状态。所以\textbf{偶数级上的动作不影响胜负}，等价 Nim 只发生在奇数级。

		\subsubsection{Lasker's Nim（可拆分的 Multi-SG）}

		\textbf{规则}：$n$ 堆，每回合二选一——\textbf{(a)} 从某堆取若干颗（标准 Nim 操作）；\textbf{(b)} 把某堆 $\ge 2$ 颗的拆成两堆非空。取到最后一颗者胜。

		\textbf{结论}：单堆 SG 函数为
		\[
			\operatorname{SG}(x) = \begin{cases}
									   x - 1, & x \bmod 4 = 0 \\
									   x,     & x \bmod 4 \in \{1, 2\} \\
									   x + 1, & x \bmod 4 = 3
			\end{cases}
		\]
		多堆答案 $= \bigoplus_i \operatorname{SG}(a_i)$，非零时先手胜。

		\textbf{打表速记}：$x = 0,1,2,3,4,5,6,7$ 时 SG 为 $0,1,2,4,3,5,6,8$——\textbf{每 $4$ 个一组、组内除了 $4k$ 与 $4k+3$ 互换以外都是恒等}。背不住时直接用通用 SG 暴力打表（见下节）现场跑。

		\subsection{巴什博弈（单堆有上限取）}

		\textbf{规则}：单堆 $N$ 颗石子，每回合取 $1 \sim M$ 颗，\textbf{取到最后一颗者胜}。

		\textbf{结论}：先手败 $\iff N \bmod (M+1) = 0$。

		\textbf{为什么是 $M+1$ 周期——"凑齐 $M+1$" 的控制术}：双方可以\textbf{稳定地}把每一回合的总取数控制为 $M+1$（一方取 $X \in [1, M]$，另一方立即取 $M+1-X$，二者必然合法）。所以 \textcolor{red}{\textbf{$N$ 被 $M+1$ 整除的局面是 P 局面}}：当前人无论怎么取，对手都能"补齐成 $M+1$"，让 $N$ 严格按 $M+1$ 递减直到 $0$（最后一颗永远被对手拿走）。$N$ 不被整除时，先手取走余数 $R = N \bmod (M+1)$ 把局面调成整除留给对手，化身"控盘者"反将一军。

		\textbf{扩展巴什博弈}：每回合取走 $X \in [a, b]$ 颗，剩余不足 $a$ 时不能再取，\textbf{取到最后一颗者胜}（注意：剩余 $R < a$ 时游戏自然结束，"最后一颗"指 $R$ 中是否含取数边界——按惯例，无法再取的当前玩家算作"未取到最后一颗"，即输家）。

		\textbf{结论}：以 $a+b$ 为周期划分，先手败 $\iff N \bmod (a+b) < a$。

		\textbf{$a+b$ 周期的直觉}：把上面"凑齐 $M+1$"换成"凑齐 $a+b$"——一方取 $X \in [a, b]$，对手取 $a+b-X \in [a, b]$ 始终合法。余数 $R \in [0, a)$ 是 P 局面：先手取走 $X \ge a$ 后剩下 $N - X$，对手补成 $a+b$ 周期把局面继续锁在余数 $<a$ 区段，最终把 $<a$ 的"垃圾尾巴"留给先手取不动。余数 $R \in [a, b]$ 时先手直取 $R$；余数 $R \in (b, a+b)$ 时先手取 $X = R - (a-1) \in (b-a+1, b]$（边界容易验合法），把局面调到余数 $a-1$ 留给对手。

		\subsection{SG 函数的暴力打表与 Anti / Every 变体}

		\textbf{通用 SG 暴力}（适合状态数 $\le 10^6$ 的题目，没法找闭式时打表观察规律）：

		\begin{minted}{cpp}
// 单堆从 x 颗一步可达的下一状态集合由 trans(x) 给出
// SG(x) = mex{ SG(y) : y in trans(x) }
int sg[N];
void compute_sg(int upTo) {
    for (int x = 0; x <= upTo; ++x) {
        vector<int> nxt = trans(x);
        vector<bool> seen(nxt.size() + 1, false);
        for (int y : nxt) if (sg[y] < (int) seen.size()) seen[sg[y]] = true;
        int m = 0; while (m < (int) seen.size() && seen[m]) ++m;
        sg[x] = m;
    }
}
		\end{minted}

		\textbf{用法}：把所有子游戏的 SG 异或起来；非零先手胜。\textbf{打表观察规律}是博弈题最常用的找通项手段——先暴力跑出 $\operatorname{SG}(0\dots 100)$ 看是否周期 / 等于 $x$ / 等于 $\lfloor x/k \rfloor$ 之类的简单形式，再去证。

		\subsubsection{反 SG（Anti-SG）}

		\textbf{规则}：与 SG 游戏一致，但\textbf{无法行动者胜}（最先不能动者赢）。

		\textbf{结论}（与反 Nim 同形）：先手胜 $\iff$ 满足其一——
		\begin{itemize}
			\item 所有子游戏 $\operatorname{SG} \le 1$ 且总 XOR $=0$；
			\item 至少一个子游戏 $\operatorname{SG} > 1$ 且总 XOR $\neq 0$。
		\end{itemize}

		\subsubsection{Every-SG（每回合走"所有还能动的石子"）}

		\textbf{规则}：DAG 上若干石子，每回合\textbf{必须把所有还能动的石子各走一步}（不能选择不动），无法行动者输。

		\textbf{结论}：对每个起点求 $\operatorname{step}(v)$——\textbf{在该子游戏中博弈能持续的回合数}（必胜方尽量缩短、必败方尽量拖长）：
		\begin{itemize}
			\item 终止节点：$\operatorname{step} = 0$。
			\item $\operatorname{SG}(v) > 0$（必胜）：$\operatorname{step}(v) = 1 + \min\{\operatorname{step}(w) : v \to w,\, \operatorname{SG}(w) = 0\}$。
			\item $\operatorname{SG}(v) = 0$（必败）：$\operatorname{step}(v) = 1 + \max\{\operatorname{step}(w) : v \to w\}$。
		\end{itemize}
		\textcolor{red}{\textbf{先手胜 $\iff$ 各子游戏 $\operatorname{step}$ 的最大值是奇数}}。直觉：所有子游戏\textbf{并行推进}，最长那局决定整盘何时结束；该局轮到谁动谁就输，奇偶反转后对应先手胜。

		\subsubsection{实战范例：CCPC 2025 网络赛 F「连线博弈」——XOR 哈希切子游戏}

		\textbf{题意}（\href{https://qoj.ac/contest/2534}{第十一届 CCPC 网络预选赛 F}）：单位圆周上等距 $n$ 个点（$2 \le n \le 10^9$），初始有 $m$ 条两两\textbf{严格不交且不共端点}的弦（$0 \le m \le 10^5$，$\sum m \le 10^5$）。Alice/Bob 轮流操作，每回合选\textbf{两个未被任何弦端点占用}的"自由点"画一条不与已有弦相交的新弦，无法操作者输。问 Alice 是否必胜。

		\textbf{子游戏分解}：每条弦把圆盘切成两半，自由点被弦端点序列划分成若干\textbf{独立子游戏}。$x$ 个自由点的子游戏一步操作占掉 $2$ 个点，并把段切成大小为 $a, b$（$a + b = x - 2$）的两个新子游戏。故
		\[
			\operatorname{SG}(x) = \mathop{\mathrm{mex}}_{a + b = x - 2} \bigl(\operatorname{SG}(a) \oplus \operatorname{SG}(b)\bigr),\quad \operatorname{SG}(0) = \operatorname{SG}(1) = 0.
		\]
		暴力打表至 $500$ 后发现 $x \ge 36$ 起呈周期 $34$，但 \textcolor{red}{\textbf{$\operatorname{idx} = 52$ 是异常项}}（$\operatorname{sg}[52] = 2$，但 $x \ge 70$ 落到 $\operatorname{idx} = 52$ 的真实周期值是 $\operatorname{sg}[86] = 9$），\textbf{必须特判}，否则会被 hack 数据 $n=8,\ m=2,\ (0,1)(4,5)$ 卡掉（朴素周期答 NO，正解 YES）。\textbf{教训}：找到看似规整的周期\textbf{务必本地用暴力对拍校验首尾几百项}，不要看一眼"前几行循环对了"就上交。

		\textbf{XOR 哈希切子游戏（关键技巧）}：题面要求"识别哪些自由点共属一个子游戏"，传统做法（按弦端点排序，相邻端点之间一段算一个子游戏）会把"两条互不相交的弦同时切到的中间区段"误判为多个独立段。\textcolor{red}{\textbf{用随机 XOR 标记代替几何判定}}——给每条弦 $(x, y)$ 生成随机权 $r$，记 $\operatorname{mp}[x] \mathbin{\hat{}{=}}\, r$、$\operatorname{mp}[y] \mathbin{\hat{}{=}}\, r$。按下标升序扫描端点，维护前缀异或 $\operatorname{cur}$；\textbf{$\operatorname{cur}$ 相同的下标段就属于同一子游戏}（因为弦的两个端点同时翻转 $\operatorname{cur}$ 两次，段内自由点的 $\operatorname{cur}$ 等于"包住它的所有未闭合弦"的随机权异或，几何上"被同一组弦包裹"等价于"在同一子游戏"）。最后对每个 $\operatorname{cur}$ 桶累加自由点数，再异或所有 $\operatorname{SG}(\text{桶大小})$。

		\begin{minted}{cpp}
// 预处理 SG 至 500
sg[0] = sg[1] = 0;
for (int i = 2; i <= 500; ++i) {
    set<int> s;
    for (int j = 0; j <= i - 2; ++j) s.insert(sg[j] ^ sg[i - 2 - j]);
    int mex = 0; while (s.count(mex)) ++mex;
    sg[i] = mex;
}

// 周期 34（从 i=36 起），idx==52 异常项特判
auto SG = [&](ll x) -> int {
    if (x <= 200) return sg[x];
    int idx = 36 + (x - 36) % 34;
    return idx == 52 ? sg[idx + 34] : sg[idx];
};

// 单组询问：XOR 哈希 + 扫描分桶
void solve() {
    cin >> n >> m;
    map<ll, ll> mp;          // 端点 → 该点累计的随机权异或
    mt19937_64 rng(time(0));
    for (int i = 0; i < m; ++i) {
        ll x, y; cin >> x >> y;
        ll r = rng();
        mp[x] ^= r; mp[y] ^= r;   // 弦两端各异或同一个 r
    }
    if (m == 0) { cout << (SG(n) ? "YES\n" : "NO\n"); return; }

    map<ll, ll> bucket;      // cur → 该组子游戏的自由点总数
    ll up = -1, cur = 0;
    mp[n] = 0;               // 末尾哨兵：把最大端点之后到 n-1 的尾段也结算
    for (auto& [u, v] : mp) {
        bucket[cur] += max(0LL, u - up - 1);  // [up+1, u-1] 的自由点数
        up = u; cur ^= v;
    }
    ll ans = 0;
    for (auto& [_, sz] : bucket) ans ^= SG(sz);
    cout << (ans ? "YES\n" : "NO\n");
}
		\end{minted}

		\textbf{为什么 XOR 哈希这一招通用}：本质是\textbf{用差分维护"被哪些区间覆盖"} —— 每条弦 $(x, y)$ 的"覆盖范围"在排序后的端点序列上是一段，而 \texttt{mp[x] \^{}= r; mp[y] \^{}= r;} 正是 \textcolor{blue}{\textbf{异或差分}}（$\operatorname{cur}$ 在 $x$ 之后翻入 $r$，在 $y$ 之后翻出 $r$，$[x, y)$ 段的 $\operatorname{cur}$ 多含 $r$）。任何"按区间覆盖签名分组"的题都可套用这套：把"判断两个位置是否属于同一组"翻译成"它们的 cur 哈希是否相等"，免掉显式区间扫线 / 并查集合并。\textbf{相关题}：HDU/CodeForces 上"求被不同区间集合覆盖的位置数 / 等价类数"系列均可一键套用。

		\subsection{威佐夫博弈（双堆联动）}

		\textbf{规则}：两堆石子 $(a, b)$（设 $a \le b$）。每回合二选一——\textbf{(a)} 从单堆取任意正数颗；\textbf{(b)} 从两堆同时取走\textbf{相同}的正数颗。取到最后一颗者胜。

		\textbf{结论}：奇异（P）局面是 $(a_k, b_k) = (\lfloor k\varphi \rfloor,\, \lfloor k\varphi \rfloor + k)$，$k = 0, 1, 2, \dots$，其中 $\varphi = \frac{1+\sqrt{5}}{2}$ 是黄金分割比。即 $(0,0), (1,2), (3,5), (4,7), (6,10), (8,13), \dots$。\textcolor{red}{\textbf{当前 $(a,b)$ 是 P 局面 $\iff$ $\lfloor (b - a) \cdot \varphi \rfloor = a$}}。

		\textbf{Beatty 定理背景}：两序列 $\{\lfloor k\varphi \rfloor\}$ 与 $\{\lfloor k\varphi^2 \rfloor\}$ 由 Beatty 定理给出"$\mathbb{N}^+$ 的不交划分"，恰好把每个正整数唯一指派为某个 $a_k$ 或 $b_k = a_k + k$。这保证了奇异局面集合"覆盖且只覆盖一次"。

		\textbf{大数精度警告}：$N \le 10^9$ 时 \texttt{double} 误差刚好够呛，建议直接用高精度常数 $\varphi \approx 1.6180339887498948482$，或直接打表黄金比的两个 Beatty 序列（项数 $\le 10^9 / \varphi \approx 6 \times 10^8$ 不行，所以现场公式判定）：

		\begin{minted}{cpp}
const double PHI = 1.6180339887498948482L;
bool first_loses(long long a, long long b) {
    if (a > b) swap(a, b);
    return (long long)((b - a) * PHI) == a;
}
		\end{minted}

		\subsection{斐波那契博弈（前手两倍约束）}

		\textbf{规则}：单堆 $N$ 颗。先手第 $1$ 回合可取 $1 \sim N-1$ 颗（不能一次拿光），此后每回合的取数上限为\textbf{上一回合所取数的 $2$ 倍}。取到最后一颗者胜。

		\textbf{结论}：\textcolor{red}{\textbf{先手败 $\iff N$ 是斐波那契数}}（$N \in \{1, 2, 3, 5, 8, 13, 21, \dots\}$）。

		\textbf{直觉——齐肯多夫（Zeckendorf）分解}：每个非斐波那契数都能唯一拆成\textbf{若干不相邻}斐波那契数之和 $N = F_{k_1} + F_{k_2} + \cdots + F_{k_t}$（$k_1 > k_2 > \dots > k_t$，$k_i - k_{i+1} \ge 2$）。先手取走最小项 $F_{k_t}$，由 Fibonacci 递推 $F_{k_{t-1}} > 2 F_{k_t}$，对手无法一次拿完下一项 $F_{k_{t-1}}$ 段，必然被迫"破坏"它，先手再取该段最小项继续逼迫……层层剥到只剩 $F_{k_1}$，此时正好对手取一颗、先手取一颗交替到底。$N$ 本身是斐波那契数时分解只有一项，先手第一手就被"取数 $\le N-1$ 不能一次拿完 + 任取必让对手反吃 $\le 2X$ 区间"的双重约束封死。

		\begin{minted}{cpp}
bool first_loses_fib(long long n) {
    long long a = 1, b = 2;
    while (b <= n) { long long c = a + b; a = b; b = c; }
    // 此时 a 是 <= n 的最大 Fibonacci
    return a == n;
}
		\end{minted}

		\subsection{图上删边游戏}

		\subsubsection{树上删边（Colon Principle）}

		\textbf{规则}：有根树，每回合删一条边，与根失去连接的部分整体被删除（即删去子树）。\textbf{无边可删者输}。

		\textbf{SG 公式}（Colon Principle，Smith 定理）：以节点 $v$ 为根的子树 SG 值 $g(v)$ 满足
		\[
			g(\text{叶子}) = 0,\qquad g(v) = \bigoplus_{c \text{ 是 } v \text{ 的孩子}} \bigl(g(c) + 1\bigr).
		\]
		注意 $+1$ 是\textbf{普通整数加法}（含义："孩子子树的 SG 等价为一堆 $g(c)$ 颗石子，挂上连父边后再加一颗"），异或才是组合多个独立分支。

		\textcolor{red}{\textbf{先手胜 $\iff g(\text{根}) \neq 0$}}（与 SG 通则一致）。

		\begin{minted}{cpp}
// 返回以 u 为根的子树 SG 值
function<int(int, int)> dfs = [&](int u, int fa) -> int {
    int s = 0;
    for (int v : adj[u]) if (v != fa) s ^= (dfs(v, u) + 1);
    return s;
};
cout << (dfs(root, -1) != 0 ? "Win\n" : "Lose\n");
		\end{minted}

		\subsubsection{无向图删边（Fusion Principle）}

		\textbf{规则}：连通无向图（可含环、重边、自环），指定一个根，每回合删一条边，与根失联的部分整体删除。无边可删者输。

		\textbf{Fusion Principle}（环约简 → 化为树上删边）：
		\begin{itemize}
			\item \textbf{奇环}（含奇数条边的环，含自环）：缩成\textbf{一个新点 + 一条挂在它上面的悬空边}。
			\item \textbf{偶环}：缩成\textbf{一个新点}（不留多余边）。
			\item 原本连接到环上任意点的所有边，全部接到这个新缩出来的点上。
		\end{itemize}

		反复约简到没有环为止，剩下的就是一棵树，套用上一节的 Colon Principle 求 SG。

		\textbf{直觉}：环上任意一条边被删后剩下一条路径——"环作为整体"的 SG 等于"等价路径的 SG"，奇数 / 偶数环的等价物正好是"$1$ 条边" / "$0$ 条边"（这是 Hackenbush 经典结论）。所以缩点保留这条等价残留即可。

	\end{multicols*}

